\documentclass[12pt,a4paper,UTF8]{article}
\usepackage{ctex}
\usepackage{geometry}
\usepackage{setspace}
\usepackage{graphicx}
\usepackage{titlesec}
\usepackage{indentfirst}
\usepackage{fancyhdr}
\usepackage{xcolor}

% ========== 页面设置 ==========
\geometry{left=2.5cm,right=2.5cm,top=2.5cm,bottom=2.5cm}
\linespread{1.5}
\setlength{\parindent}{2em}

% ========== 修复 fancyhdr 报错 ==========
\setlength{\headheight}{14pt}

% ========== 页眉 ==========
\pagestyle{fancy}
\fancyhf{}
\lhead{\small 哈尔滨工业大学（深圳）马克思主义学院课程论文}
\rhead{}
\cfoot{\thepage}
\renewcommand{\headrulewidth}{0.4pt}

% ========== 标题格式 ==========
\titleformat{\section}{\centering\bfseries\zihao{3}}{}{0em}{}
\titleformat{\subsection}{\bfseries\zihao{4}}{}{0em}{}

\begin{document}
\pagestyle{empty}

% ====================== 封面 ======================
\begin{titlepage}
    \centering
    \includegraphics[width=0.5\textwidth]{校徽.png}\\[0.8cm]
    {\Large 马克思主义学院}\\[1cm]
    {\Large 2026年春季学期}\\[0.3cm]
    {\Large 本科生课程考核}\\[0.3cm]
    {\Large （课程论文）}\\[2cm]

    \begin{tabular}{l l}
        课程名称： & 形势与政策 \\
        论文题目： & 《端牢能源饭碗：新型能源体系下新旧能源协同发展》 \\
        学生所在院（系）： &  \\
        专业： &  \\
        学号： &  \\
        姓名： &  \\
        上交日期： &  \\
    \end{tabular}
    \vfill
\end{titlepage}

\pagestyle{fancy}

% 醒目防抄袭声明
\noindent{\bfseries\color{red}重要声明：本文仅作为课程写作参考资料，严禁抄袭、照搬用于课程提交，抄袭后果自负！}\\[1em]

% ====================== 摘要 ======================
\newpage
\begin{center}
    \bfseries\zihao{3} 摘要
\end{center}
\vspace{0.5cm}
\zihao{4}
\indent 二十届四中全会将统筹发展和安全、健全新型能源体系纳入国家战略重点，2026年全国两会再次着重强调要扎实保障能源安全、更好端牢能源饭碗，推动新能源与传统能源统筹协调、协同发展。国家安全以人民安全为宗旨，生命安全、民生保障是一切发展的根本前提，能源安全作为非传统安全的重要组成，直接关系社会稳定、民生福祉与人民生命财产安全。当前世界百年未有之大变局加速演进，中东地缘局势持续紧张，2026年美伊冲突不断升级，霍尔木兹海峡作为全球能源运输要道面临通航风险，国际油价震荡波动，全球能源供给格局受到明显冲击，进一步影响民生保障与社会安全稳定。与此同时，美国前总统特朗普拟于5月访华，中美在经贸、科技、地区安全以及能源战略层面的互动博弈，深刻影响我国能源发展外部环境与国家安全大局。

本文立足形势与政策课程四大学习专题，以专题四“更好端牢能源的饭碗”为核心，结合总体国家安全观中生命安全、人民安全核心要义，围绕新型能源体系建设、新能源与传统能源协同发展展开研究，结合二十届四中全会精神与2026年全国两会部署，剖析美伊冲突升级、特朗普访华带来的外部机遇与挑战，从能源供给多元化、新旧能源协同转型、科技自立自强、中东外交布局、中美关系稳健把控等方面提出实践路径，坚守人民至上、生命至上的安全底线，为理解我国能源安全战略、构建安全高效的新型能源体系、筑牢民生与生命安全屏障提供参考。

\vspace{0.5cm}
\noindent\textbf{关键词：} 能源饭碗；新型能源体系；新能源；传统能源协同；美伊冲突；中美关系；总体国家安全观；生命安全

% ====================== 正文 ======================
\newpage
\zihao{4}

\section{引言}
二十届四中全会明确提出要完善国家安全体系、健全新型能源体系，把能源安全作为统筹发展和安全的重要抓手，为我国能源长远发展划定顶层设计。总体国家安全观明确坚持以人民安全为宗旨，人民生命安全、生活保障、民生稳定是国家安全的核心落脚点，能源稳定供给是保障社会正常运转、守护人民生命财产安全的基础保障。2026年全国两会政府工作报告进一步落实部署，强调必须牢牢守住能源安全底线，坚持在发展中保障安全，推动风电、光伏、水电等新能源规模化发展，同时做好传统能源清洁高效利用，实现新能源与传统能源互补协调，真正把能源饭碗牢牢端在自己手中，以能源安全护航民生安全、生命安全与社会稳定。

当前国际地缘格局深刻调整，中东地区凭借关键地理位置与富集油气资源，长期成为大国博弈的核心区域。2026年美伊矛盾持续激化，军事摩擦频发，霍尔木兹海峡航运安全承压，国际油价起伏不定，给全球产业链、供应链稳定带来新的不确定性，间接影响国内物资供给、民生保障与公共安全。与此同时，美国前总统特朗普即将访华，中美在经贸合作、科技竞争、台湾问题、中东安全以及全球能源治理等领域的分歧与交流，直接塑造我国能源发展的外部环境与国家安全态势。

中东变局与中美互动并非孤立的国际事件，而是与我国新型能源体系建设、新旧能源协同发展、人民生命民生保障高度关联。我国作为全球最大能源进口国和新能源产业大国，中东局势直接影响化石能源进口安全，进而影响工业生产、民生用电用气保障，关系人民基本生活与社会公共安全；中美战略博弈则影响能源技术创新、产业链自主可控与国际能源合作空间。基于形势与政策课程学习要求，紧扣“更好端牢能源的饭碗”专题，结合总体国家安全观生命安全核心要义、二十届四中全会与2026年全国两会精神，以新能源与传统能源协同发展为切入点，分析外部局势带来的机遇挑战，并提出对应应对策略，具有重要现实意义与学习价值。

\section{美伊冲突升级对我国能源安全与发展的影响}
\subsection{能源对外依存度偏高，倒逼新型能源体系加速建设}
能源是国民经济的命脉，是保障民生运转、守护社会稳定与人民生命财产安全的基础支撑，保障能源安全是二十届四中全会和2026年全国两会重点强调的战略任务，更是落实总体国家安全观、践行人民至上、生命至上理念的关键举措。我国原油、天然气对外依存度居高不下，中东是我国化石能源进口最主要来源地，霍尔木兹海峡承担全球大量石油运输任务。美伊冲突升级导致地区局势紧张，航线通行风险加大、运输成本攀升，国际油价走高，直接推高国内交通运输、化工制造等行业生产成本，影响民生物资稳定供给，给社会稳定与民生安全带来潜在隐患。

外部地缘风险的不确定性，倒逼我国加快构建新型能源体系，践行“新能源扩容、传统能源提质”的发展思路，不再过度依赖中东单一化石能源供给，通过新旧能源协同发展夯实能源安全根基，以稳定能源供给守住民生底线、守护人民生活与生命安全保障。

\subsection{一带一路中东合作机遇与风险并存}
中东是共建“一带一路”的重要枢纽区域，我国与中东国家在能源贸易、基建投资、产能合作等领域合作基础扎实。美伊局势动荡造成部分区域政局不稳、安全风险上升，一定程度上延缓合作项目落地进度，增加海外投资与运营风险。

但从长远来看，地区动荡也促使中东各国更加重视经济转型与能源结构优化，主动寻求与我国在新能源发电、电力装备、绿色基建等领域合作。我国坚持不干涉别国内政的外交原则，秉持客观公正立场参与中东事务，契合二十届四中全会提出的高水平对外开放要求，为我国输出新能源技术、拓展国际能源合作创造良好契机，助力国内外能源安全体系协同建设，全方位守护发展安全与民生安全。

\subsection{地缘博弈复杂化，要求提升能源外交战略话语权}
中东地区民族宗教矛盾交织、大国博弈错综复杂，美伊冲突背后既有双边历史恩怨，也有美国掌控中东秩序、遏制区域发展的战略意图。我国在中东外交中需要平衡与伊朗、沙特、美国等多方关系，坚持不选边、不站队，维护自身能源利益与战略空间。

同时积极参与中东热点问题斡旋，倡导通过对话协商化解分歧，践行共同、综合、合作、可持续的安全观，在全球能源治理中发出中国声音，为新型能源体系建设营造稳定的地缘政治环境，以外部安全保障内部发展安全、民生安全与生命安全。

\section{特朗普5月访华对我国能源战略布局的影响}
\subsection{经贸与能源议题交织，双边合作与摩擦并存}
中美经贸关系是两国关系的压舱石，能源贸易长期是中美合作的重要组成部分。特朗普此次访华，经贸、市场准入、产业规则、能源合作等议题将成为磋商重点。美方可能在新能源产业规则、贸易壁垒等方面对华施压，限制我国新能源产业国际化发展空间。

与此同时，中美在能源贸易、气候变化、绿色低碳发展等领域存在共同利益，具备合作基础。会晤有利于管控分歧、稳定双边关系，为我国新能源产业出海、能源装备贸易、参与全球能源治理营造相对平稳的外部环境，契合2026年全国两会扩大高水平对外开放的工作部署，助力我国以能源安全筑牢民生与生命安全防线。

\subsection{科技竞争延伸至能源领域，倒逼能源科技自立自强}
中美科技博弈已延伸到新能源材料、集成电路、高端电机电控、储能技术、电力电子等能源核心领域。美国通过技术管制、产业链封锁等方式，试图遏制我国新能源产业升级步伐。

二十届四中全会明确要求强化国家战略科技力量，加快关键核心技术攻关。面对外部遏制，我国必须坚持创新驱动，突破新能源发电、智能电网、大容量储能等关键技术，以科技自立自强支撑新型能源体系建设，依托技术自主可控保障能源稳定供给，从根源上守护人民生活保障与社会安全稳定。

\subsection{中东热点成博弈交汇点，中美分歧与协调空间同在}
美伊冲突、霍尔木兹海峡安全、伊朗核问题等中东事务，将成为中美高层会晤绕不开的议题。美方坚持单边制裁与军事施压，中方主张停火止战、政治解决争端，双方立场存在明显分歧。但在保障海峡通航安全、稳定国际能源市场、打击恐怖主义等方面，中美拥有共同利益，具备沟通协调空间。

我国将继续秉持客观公正立场，推动中东局势降温，为国内能源进口安全、新旧能源协同发展争取稳定外部条件，持续夯实国家安全与民生生命安全保障基础。

\section{立足能源饭碗视角的系统性应对策略}
\subsection{统筹新旧能源布局，推进新能源与传统能源协同发展}
紧扣2026年全国两会“端牢能源饭碗”核心要求，落实二十届四中全会健全新型能源体系的部署，坚持新能源和传统能源两条腿走路，立足总体国家安全观人民至上、生命至上核心宗旨，以能源安全保障民生安全、社会安全。一方面大力发展风电、光伏、水电、核电及储能产业，扩大新能源装机规模，提升新能源消纳能力；另一方面做好煤炭、油气资源清洁高效利用，推进传统能源绿色转型。通过新旧能源互补调峰、协同保供，优化国家能源结构，从根本上降低地缘局势波动带来的供给风险，保障全社会能源稳定供应，守护人民正常生产生活秩序。

\subsection{拓宽多元进口渠道，分散中东地缘政治风险}
针对美伊冲突带来的航运与供给不确定性，持续推进能源进口来源、运输通道多元化。稳固中东传统能源合作，同时加大俄罗斯、中亚、非洲、拉美等地区能源进口布局，完善陆海双通道运输网络；深化海外能源产能合作与股权投资，提升能源供给韧性，把能源安全主动权牢牢掌握在自己手中，筑牢民生保障与生命安全的能源根基。

\subsection{坚持创新驱动，以科技支撑新型能源体系建设}
对标二十届四中全会创新驱动战略，聚焦新能源发电、智能电网、高端电气装备、储能技术等关键领域加大科研投入，推动产学研深度融合。深耕能源核心技术攻关，突破“卡脖子”难题，以能源科技自主可控保障能源体系安全稳定，为守护社会稳定、人民生命财产安全提供硬核科技支撑。

\subsection{深化中东务实合作，拓展能源外交新格局}
持续深化与中东国家全方位务实合作，在传统油气合作基础上，拓展风电光伏、电力基建、绿色能源项目合作，输出中国技术与中国方案。坚持不干涉内政原则，积极参与中东和平进程，以能源合作为纽带增进政治互信，为我国新型能源体系对外开放营造良好地缘环境，以外部安全护航内部民生安全与生命安全保障。

\subsection{稳健把控中美关系，营造能源转型稳定外部环境}
秉持相互尊重、和平共处、合作共赢原则，通过高层常态化沟通管控经贸、科技、地区热点分歧。在气候治理、绿色低碳、全球能源治理等领域寻求合作共识，同时坚持底线思维，做强自身能源产业与科技实力，不受外部博弈牵制，保障国内能源转型长期平稳推进，持续夯实国家安全与人民生命安全保障体系。

\section{结论}
二十届四中全会为我国新型能源体系建设、统筹发展与安全指明战略方向，2026年全国两会进一步细化部署，强调要始终端牢能源饭碗，推动新能源与传统能源协同有序发展。总体国家安全观以人民安全为宗旨，生命安全、民生稳定是国家安全的根本落脚点，能源安全作为重要的非传统安全领域，直接关乎社会运转、民生保障与人民生命财产安全。2026年美伊冲突升级、特朗普拟访华两大国际事件相互交织，既给我国能源供给安全、产业发展、外交布局、民生生命安全保障带来外部挑战，也为我国加快能源结构转型、拓展国际能源合作、完善国家安全体系提供战略机遇。

面对百年变局与地区动荡叠加的复杂形势，我国必须坚守能源安全底线，以构建新型能源体系为核心，坚持新能源扩容提质与传统能源清洁利用协同并进，拓宽多元能源供给渠道，强化能源科技自立自强，深化中东能源务实合作，稳健应对中美战略博弈，持续提升国家能源安全保障能力，以稳固的能源安全守护社会稳定、民生幸福与人民生命安全。

新时代青年应深入学习领会二十届四中全会与全国两会精神，深刻把握总体国家安全观内涵，树立正确能源观与国际大局观，心系国家能源安全与民生安全建设，主动担当时代使命，助力能源强国与安全中国建设。

% ====================== 参考文献 ======================
\newpage
\begin{center}
    \bfseries\zihao{3} 参考文献
\end{center}
\vspace{0.5cm}
\zihao{4}
\begin{enumerate}
    \item 习近平. 论坚持推动构建人类命运共同体[M]. 北京：中央文献出版社，2023.
    \item 中共中央. 中国共产党第二十届中央委员会第四次全体会议公报[R]. 2025.
    \item 国务院. 2026年全国两会政府工作报告[R]. 2026.
    \item 国家能源局. 加快构建新型能源体系 牢牢端牢能源饭碗[N]. 中国能源报，2026.
    \item 外交部. 中美双方就特朗普总统访华事保持沟通[EB/OL]. 2026-05-06.
    \item 央视网. 今日关注：中东战事升级与海峡对峙[EB/OL]. 2026-04-18.
    \item 环球时报. 美伊冲突升级影响全球能源格局[N]. 环球时报，2026-04-20.
\end{enumerate}

\end{document}